\setauthor{Markus Tran}
\section{Datenbanken}
Um Dateneingaben persistent zu speichern und strukturiert zu verarbeiten, ist der
Einsatz einer Datenbank unerlässlich. Eine Datenbank stellt die zentrale Instanz
für die Speicherung, Organisation und Verwaltung von Daten dar. Sie ermöglicht
nicht nur das einfache Ablegen von Informationen, sondern auch deren gezielte
Abfrage, Aktualisierung und Löschung. Um auf die gespeicherten Daten zugreifen
und diese verwalten zu können, wird ein sogenanntes Datenbankmanagementsystem
(DBMS) eingesetzt, das als Vermittler zwischen der Anwendung und den physisch
gespeicherten Daten fungiert.

\subsection{Datenbank definition}
Eine Datenbank ist eine strukturierte Sammlung von Daten, die so organisiert ist,
dass sie effizient gespeichert, abgerufen und verwaltet werden kann. Im
Gegensatz zur einfachen Dateiablage bietet eine Datenbank Mechanismen zur
Vermeidung von Redundanzen, zur Sicherstellung der Datenkonsistenz sowie zur
gleichzeitigen Nutzung durch mehrere Benutzer oder Anwendungen.

Grundlegende Begriffe im Datenbankumfeld sind:

\begin{itemize}
    \item \textbf{Entität}: Ein eindeutig identifizierbares Objekt der realen Welt,
    das in der Datenbank abgebildet wird (z.\,B. ein Benutzer oder ein
    Streaming-Anbieter).
    \item \textbf{Attribut}: Eine Eigenschaft einer Entität (z.\,B. Name oder ID).
    \item \textbf{Relation}: Eine Tabelle in einer relationalen Datenbank, die
    Entitäten eines bestimmten Typs speichert.
    \item \textbf{Primärschlüssel}: Ein Attribut oder eine Kombination von
    Attributen, das eine Entität eindeutig identifiziert.
    \item \textbf{Fremdschlüssel}: Ein Attribut, das auf den Primärschlüssel einer
    anderen Tabelle verweist und damit Beziehungen zwischen Entitäten abbildet.
\end{itemize}



\subsection{Datenbank Managment System - DBMS}

Ein Datenbankmanagementsystem (DBMS) ist eine Software, die den Zugriff auf eine
Datenbank verwaltet und kontrolliert. Es stellt die Schnittstelle zwischen der
Anwendung und den physisch gespeicherten Daten dar und übernimmt Aufgaben wie
die Speicherverwaltung, die Zugriffskontrolle, die Transaktionsverwaltung sowie
die Sicherstellung der Datenkonsistenz.

Ein DBMS besteht im Wesentlichen aus zwei zentralen Komponenten. Der
\textit{physische Datenbestand} umfasst die tatsächlich auf einem Datenträger
gespeicherten Daten. Der \textit{Systemkatalog} -- auch Data Dictionary genannt
-- enthält Metainformationen über die Struktur der Datenbank, also Informationen
über Tabellen, Spalten, Datentypen, Beziehungen und Zugriffsrechte. Mithilfe des
Systemkatalogs kann das DBMS Anfragen korrekt interpretieren und die gespeicherten
Daten strukturiert zurückgeben.

Zu den bekanntesten DBMS zählen PostgreSQL, MySQL, Oracle Database und
Microsoft SQL Server im relationalen Bereich sowie MongoDB und Redis im
nicht-relationalen Bereich.

\subsection{Datenbank Komponenten}
Die Physische Speicherung aka der Datenbaestand wo was gespeichert ist und der Systemkatalog mit dem 
ein DBMS dann auf die Datenzugreifen konnte.
